\documentclass[25pt, a0paper, portrait, margin=0mm, innermargin=15mm, blockverticalspace=10mm, colspace=10mm, subcolspace=8mm]{tikzposter}

\usepackage{amsmath}
\usepackage{epsf,graphicx,subfig}

\usepackage{amsmath}
\usepackage{epsf,graphicx,subfig}
\setcounter{tocdepth}{3}
\usepackage{xcolor}
\usepackage{lineno}
\usepackage[nolist]{acronym}
\usepackage{amssymb,amsmath}
\usepackage{standalone}
\usepackage{multicol}
\usepackage{tikz}
\usepackage{pgfplots,pgfplotstable}
\usepackage{scalefnt}
\usetikzlibrary{positioning}
\usetikzlibrary{pgfplots.groupplots}
\usepackage{url}
\pgfplotsset{compat=1.8}

\usepackage{adjustbox}
\usepackage{epstopdf}

% include pgfplots and do a workaround to add the missing layers. This problem might be related to pgfplots and tikz versions as pointed out in: http://tex.stackexchange.com/questions/104010/why-does-loading-pgfplots-after-tikz-break-the-default-layers-in-a-tikzpicture

\usepackage{pgfplots}

% define layers
    \pgfdeclarelayer{foreground}
    \pgfdeclarelayer{background}
    \pgfdeclarelayer{backgroundlayer} 
    \pgfdeclarelayer{notelayer}
% tell TikZ how to stack them (back to front)
    \pgfsetlayers{background,backgroundlayer,main,notelayer,foreground}


\definecolor{udgColor}{RGB}{82,119,213}

\title{\parbox{0.8\linewidth}{\centering \acs{sift} texture description for understanding breast ultrasound images}}
\author{\parbox{0.8\linewidth}{\centering Joan Massich\textsuperscript{1,2},
                                          Fabrice Meriaudeau\textsuperscript{2}, 
                                          Melcior Sent{\'i}s\textsuperscript{3},
                                          Sergi Ganau\textsuperscript{3},
                                          Elsa~P{\'e}rez\textsuperscript{4},
                                          Domenec Puig\textsuperscript{5},
                                          Robert  Mart{\'i}\textsuperscript{1},
                                          Arnau Oliver\textsuperscript{1}
                                          and Joan Mart{\'i}\textsuperscript{1}}}
\institute{ 
  \textsuperscript{1}Computer Vision and Robotics Group, University of Girona, Spain.
  \textsuperscript{2}Laboratoire Le2i-UMR CNRS, University of Burgundy,  Le Creusot, France.
  \textsuperscript{3}Department of Breast and Gynecological Radiology,  UDIAT-Diagnostic Center, Parc Taul{\'i} Corporation, Sabadell, Spain.
  \textsuperscript{4}Department of Radiology, Hospital Josep Trueta of Girona, Spain.
  \textsuperscript{5}Department of Computer Engineering and Mathematics, University Rovira i Virgili, Tarragona, Spain.
}
\titlegraphic{
\includegraphics[width=5cm]{figures/logos/udg.png}
 \hfill 
\includegraphics[width=5cm]{figures/logos/uB.pdf}
\\% \vspace{5pt} \\
\hfill
\\
\includegraphics[width=5cm]{figures/logos/urv.pdf}
 \hfill 
\includegraphics[width=5cm]{figures/logos/ics.png}
}


% Fix the logos positioning
\makeatletter
\def\TP@titlegraphictotitledistance{-9cm}   % -9 can be adjusted to one's need
\settitle{ \centering \vbox{
\@titlegraphic \\[\TP@titlegraphictotitledistance] \centering
\color{titlefgcolor} {\bfseries \Huge \sc \@title \par}
\vspace*{1em}
{\huge \@author \par} \vspace*{1em} {\small \@institute}
}}
\makeatother

%\institute{contact author: sik@eia.udg.edu}
%\usetheme{Autumn}\usecolorstyle[colorPalette=BrownBlueOrange]{Germany}
\usetheme{Autumn} \usecolorstyle[colorOne=udgColor,colorTwo=udgColor!15!white]{Germany} 

\input{defineBreastGTlabelColors.tex}

\begin{document}\maketitle
/home/sik/Work/escola/recerca/iwdm2014/paper/acronyms.tex
\acresetall

\graphicspath{{figures/paperFigures/}}


\begin{columns} 
  \column{0.6}
\block{Summary}{
Texture is a powerful cue for describing structures that show a high degree of similarity in their image intensity patterns. 
This work describes the use of \acf{sift}, both as low-level and high-level descriptors, applied to differentiate the tissues present in breast US images. 
For such a task, a subset of 16 images has been randomly selected from a larger dataset of 700 \ac{us} images acquired at the \emph{UDIAT Diagnostic Centre of Parc Taul\'{i}} in Sabadell (Catalunya), between 2010 and 2012.
This subset has been complemented with multi-label \ac{gt}, as illustrated in figure~\ref{fig:dataExample}.

%Experimental results are provided showing the validity of the proposed approach for describing the tissues in breast US images.

}
  \column{0.4}
\block{}{
\begin{tikzfigure}[Data sample: image, accompanying multi-label \ac{gt}, tissue label \ac{gt} color-coding.]
  \centering
  \documentclass[border=2pt]{standalone}
\usepackage{tikz}


\input{../../defineBreastGTlabelColors.tex}

\begin{document}
\begin{tikzpicture}

\node[anchor=south west,inner sep=0] (imgNode) at (0,0) {\includegraphics[width=.12\textwidth]{pngImgs/gt/000002.png}};

\node[anchor=south east,inner sep=0] at (-8pt,0) {\includegraphics[width=.12\textwidth]{pngImgs/000002.png}};

\definecolor{lungColor}{rgb}{0.2353, 0.6078, 0.8235}
\definecolor{chestWallColor}{rgb}{0.5294, 0.7843, 0.6078}
\definecolor{ribColor}{rgb}{0.0000, 0.4510, 0.1961}
\definecolor{pectoralColor}{rgb}{0.6510, 0.3490, 1.0000}
\definecolor{fibroGlandColor}{rgb}{1.0000, 1.0000, 0.0000}
\definecolor{fatColor}{rgb}{0.9804, 0.5882, 0.1176}
\definecolor{skinColor}{rgb}{0.9804, 0.7255, 0.7451}
\definecolor{unkTissueColor}{rgb}{0.6000, 0.3020, 0.2510}
\definecolor{bgColor}{rgb}{0.0000, 0.0000, 0.0000}
\definecolor{lesionColor}{rgb}{1.0000, 0.2510, 0.0000}
\definecolor{boundaryColor}{rgb}{0.8784, 0.8784, 0.7529}

\tikzset{nameSt/.style=
{anchor=north west,rectangle,
node distance=15pt,
minimum width=8pt,minimum height=8pt,
}}
%\begin{scriptsize}
%
\draw[] (imgNode.north east) +(8pt,0) node[nameSt, draw=chestWallColor, fill=chestWallColor,label=right:Chest wall] (cwName) {} ;
\draw[] node[nameSt, draw=ribColor, fill=ribColor,below=of cwName,label=right:Rib] (ribName) {} ;
\draw[] node[nameSt, draw=skinColor, fill=skinColor,below=of ribName,label=right:Skin layers] (skinName) {} ;
%\draw[] node[nameSt, draw=unkTissueColor, fill=unkTissueColor,below=of skinName,label=right:Chest wall] (unkTissueName) {} ;

\draw[] node[nameSt, draw=lesionColor, fill=lesionColor,below=of skinName,label=right:Lesion] (lesionName) {} ;
\draw[] node[nameSt, draw=boundaryColor, fill=boundaryColor,below=of lesionName,label=right:Boundary] (boundaryName) {} ;

\draw[] (imgNode.north east) +(8,0) node[nameSt,draw=lungColor, fill=lungColor,label=right:Air or lungs] (lungName) {};
\draw[] node[nameSt, draw=pectoralColor, fill=pectoralColor,below=of lungName,label=right:Pectoral muscle] (pectoralName) {} ;
\draw[] node[nameSt, draw=fibroGlandColor, fill=fibroGlandColor,below=of pectoralName,label=right:Fibro-glandular tissue] (fibroGlandName) {} ;
\draw[] node[nameSt, draw=fatColor, fill=fatColor,below=of fibroGlandName,label=right:Adipose tissue] (fatName) {} ;
\draw[] node[nameSt, draw=bgColor, fill=bgColor,below=of fatName,label=right:Background] (bgName) {} ;
%
%\end{scriptsize}

\end{tikzpicture}
\end{document}

  \label{fig:dataExample}
\end{tikzfigure}
}
\end{columns}
%\begin{keywords}
%breast cancer, ultrasound, texture, SIFT
%\end{keywords}

\block{\ac{sift} as a low-level descriptor, tested using \ac{map}}{
  \begin{multicols*}{3}
    \begin{itemize}

      \item Goal: to analyze tissue classes separability based on
        \begin{itemize}
          \item Bayesian framework for tissue discrimination (see eq.\,1).%\ref{eq:bayes}).
          \item low-level descriptors based on \ac{sift}
        \end{itemize}
    \end{itemize}

    \begin{equation*}
      P(\omega|\bar{x}) = \frac{P(\bar{x}|\omega)\cdot P(\omega)}{P(\bar{x})} \qquad \text{(1)}
      \label{eq:bayes}
    \end{equation*}

    \begin{itemize}
      \item Feature Description.
        \begin{itemize}
          \item Extract \ac{sift} descriptors at all pixel positions.
          \item Project the $128D$ \ac{sift} to $2D$ using \ac{pca}.
        \end{itemize}
        \vfill
        \columnbreak
        
      \item Qualitative analysis can be found in fig.\,\ref{fig:model},\ref{fig:prior} and \ref{fig:MAP}.
      \item Quantitative analysis can be found in fig.\,\ref{fig:llConfMatrix},
            where a comparison with intensity feature can be found. The overall sensitivity for the intensity case is $16.6\pm27.5\%$, whereas for the \ac{sift} case is $18.8\pm17.2\%$ which shows that both feature spaces produce similar results.
    
    \end{itemize}
    
  \end{multicols*}
}

\begin{columns} 

\column{0.34} \block{}{
\begin{tikzfigure}[Distribution of the \acs{sift} descriptors for some classes in the \ac{gt}.]
\documentclass[border=2pt]{standalone}
\usepackage{tikz}


\input{../../defineBreastGTlabelColors.tex}

\begin{document}
 \begin{tikzpicture}
\tikzset{myNode/.style=
{anchor=north west,rectangle,
node distance=5pt,
minimum width=8pt,minimum height=8pt,
inner sep=0,
}}

\newcommand\mySize{6cm}

\node[myNode, label=below:\small Background] (bgNode) at (0,0) {\includegraphics[width=\mySize]{gtDistro/000.png}};
\node[myNode, right=of bgNode, label=below:\small tiny Air or lungs] (airNode) {\includegraphics[width=\mySize]{gtDistro/001.png}};
\node[myNode, right=of airNode, label=below:\small Chest wall] (cwNode) {\includegraphics[width=\mySize]{gtDistro/002.png}};

\node[myNode, right=of cwNode, label=below:\small Rib] (ribNode) {\includegraphics[width=\mySize]{gtDistro/003.png}};

\node[myNode, below=40pt of bgNode, label=below:\small Fibro-glandular] (fibNode) {\includegraphics[width=\mySize]{gtDistro/005.png}};

\node[myNode, right=of fibNode, label=below:\small Adipose tissue] (fatNode) {\includegraphics[width=\mySize]{gtDistro/006.png}};
\node[myNode, right=of fatNode, label=below:\small Skin layers] (skNode) {\includegraphics[width=\mySize]{gtDistro/007.png}};
\node[myNode, right=of skNode, label=below:\small Lesion] (lesionNode) {\includegraphics[width=\mySize]{gtDistro/128.png}};

\end{tikzpicture} 
\end{document}

  \label{fig:model}
\end{tikzfigure}
}

\column{0.14} \block{}{
\begin{tikzfigure}[Data prior knowledge.] %Class tiussue and feature priors.]
  \documentclass[border=2pt]{standalone}
\usepackage{tikz}
\usepackage{pgfplots,pgfplotstable}
\pgfplotsset{compat=1.4}
\usetikzlibrary{pgfplots.groupplots}

%\graphicspath{{../paperFigures/}}
\input{../../defineBreastGTlabelColors.tex}

\begin{document}

\pgfplotstableread[col sep=comma]{
Class, Occurrence, Pw
%-------------------------------------
Bg,     518738,	0.207
Lungs,  447153,	0.178
CWall,  134165,	0.054
Ribs,   22541, 	0.009
Pect,	  182468,	0.073
FGlan,  525238,	0.210
Fat,    413118,	0.165
Skin,   17732, 	0.007
Lsion,  163848,	0.065
Bound,  81082, 	0.032
}\occurrenceDataTable

\pgfplotscreateplotcyclelist{colorbrewer-RYB}{
{bgColor!50!black,fill=bgColor},
{lungColor!50!black,fill=lungColor},
{chestWallColor!50!black,fill=chestWallColor},
{ribColor!50!black,fill=ribColor},
{pectoralColor!50!black,fill=pectoralColor},
{fibroGlandColor!50!black,fill=fibroGlandColor},
{fatColor!50!black,fill=fatColor},
{skinColor!50!black,fill=skinColor},
{lesionColor!50!black,fill=lesionColor},
{boundaryColor!50!black,fill=boundaryColor},
}

\pgfplotsset{
    select row/.style={
        x filter/.code={\ifnum\coordindex=#1\else\def\pgfmathresult{}\fi}
    }
}

\newcommand\mySize{6cm} 

\begin{tikzpicture}
\begin{axis}[ 
              ybar=0pt,
              /pgf/bar shift=0pt,scale only axis,
              width=\mySize,
              ymin=0,
              xtick=\empty,
              tickpos=left,
              scaled y ticks=base 10:2,
              cycle list name=colorbrewer-RYB,
              bar width=14pt,
              ymajorgrids,
              font=\tiny,
              title=\small $P(\omega)$,
              name=pwAxis,
            ] 


\pgfplotsinvokeforeach  {0,...,9}{
  \addplot table [ x expr=\coordindex, select row=#1, y=Pw] {\occurrenceDataTable};
}
\end{axis}

\node at (pwAxis.south) 
      [ anchor=north,
        inner sep=0,
        yshift=-10pt,
%        draw,
        label=below:\small $P(\bar{x})$,
      ]{\includegraphics[height=\mySize]{siftOccurrences2}};
\end{tikzpicture}
\end{document}

  \label{fig:prior}
\end{tikzfigure}
} 

\column{0.215} \block{}{
\begin{tikzfigure}[ Qualitative evaluation of the \ac{map} labeling of the feature space.]
\documentclass[border=2pt]{standalone}
\usepackage{tikz}

\begin{document}
\begin{tikzpicture}
\node[anchor=south west,inner sep=0] (imgNode) at (0,0) {\includegraphics[trim = 91 230 90 230, clip,width=10cm]{labelingMAP.pdf}};

%\node[anchor=south east,inner sep=0] {\includegraphics[trim = 85 230 180 230, clip,width=.1\textwidth]{intensityLabelingMAP.pdf}};
\end{tikzpicture}
\end{document}

  \label{fig:MAP}
\end{tikzfigure}
}

\column{0.305} \block{}{
\begin{tikzfigure}[Confusion matrix showing quantitative results obtained from ($10.000\text{ samples} \times 10 \text{ classes}) \times 5\text{-fold}$ cross-validation. 
]
\documentclass[border=2pt]{standalone}
\usepackage{tikz}
\usepackage{pgfplots}

\definecolor{bgColor}{rgb}{0.0000, 0.0000, 0.0000}
\definecolor{boundaryColor}{rgb}{0.8784, 0.8784, 0.7529}
\definecolor{chestWallColor}{rgb}{0.5294, 0.7843, 0.6078}
\definecolor{fatColor}{rgb}{0.9804, 0.5882, 0.1176}
\definecolor{fibroGlandColor}{rgb}{1.0000, 1.0000, 0.0000}
\definecolor{lesionColor}{rgb}{1.0000, 0.2510, 0.0000}
\definecolor{lungColor}{rgb}{0.2353, 0.6078, 0.8235}
\definecolor{pectoralColor}{rgb}{0.6510, 0.3490, 1.0000}
\definecolor{ribColor}{rgb}{0.0000, 0.4510, 0.1961}
\definecolor{skinColor}{rgb}{0.9804, 0.7255, 0.7451}
\definecolor{unkTissueColor}{rgb}{0.6000, 0.3020, 0.2510}

\begin{document}
\begin{tikzpicture}
\newcommand\confMatrixSize{7cm}
\newcommand\gtLabelDistance{0.3cm}

\tikzset{
every pin/.style={fill=yellow!50!white,rectangle,rounded corners=3pt,font=\tiny},
small dot/.style={fill=black,circle,scale=0.3}
}

\pgfplotsset{
	myStyleCommon/.append style={	width=\confMatrixSize,height=\confMatrixSize,
									scale only axis,
									xlabel={Estimated Class},
								    ylabel={Actual Class},
								    yticklabels={},
								    xticklabels={},
									xticklabel pos=right,
									xlabel near ticks,
									xtick={1.5,2.5,...,9.5},
									ytick={1.5,2.5,...,9.5},
									grid,
      }
}

\pgfplotsset{
	confMatrix/.append	style={	myStyleCommon,
								axis on top,
								colormap/jet,
								point meta min=0,
								point meta max=1,
								xmin=0.5,
								xmax=10.5,
								y dir=reverse,
								ymin=0.5,
								ymax=10.5,
							  }
}
\pgfplotsset{							  
	gtLabelsAxis/.append style={myStyleCommon,
%								axis lines=none,
								colormap={labelsGT}{color(0cm)=(bgColor);
					 								color(1cm)=(lungColor);	
					 								color(2cm)=(chestWallColor);
										 			color(3cm)=(ribColor);
										 			color(4cm)=(pectoralColor);
										 			color(5cm)=(fibroGlandColor); 
										 			color(6cm)=(fatColor); 
					 								color(7cm)=(skinColor);    
								 					color(8cm)=(lesionColor);
							 						color(9cm)=(boundaryColor); 
					 								color(10cm)=(unkTissueColor); 
					 								},
							    point meta min=0,
							    point meta max=10,
							    xmin=-0.5,xmax=9.5,
%							    height=0cm,
							  }				  
}
\tikzstyle{gtLabelsPlotStyle}=[scatter,only marks,mark=square*,mark size=5pt,domain=0:9,samples at={0,...,9}]


\begin{axis}[confMatrix,name=siftConfMatrix,clip = false,
every axis x label/.style= {at={(axis description cs:0.5,1.1)},draw=red,anchor=north},
every axis y label/.append style= {at={(axis description cs:-0.1,0.5)},draw=red,anchor=east},
yticklabel style={draw=red},
xticklabel style={draw=red},
]
\addplot graphics [xmin=0.5,xmax=10.5,ymin=0.5,ymax=10.5] {LowLevelSIFTConfusionMatrixMean.png};

\addplot[
scatter,
only marks,
mark size=3pt,
scatter src=explicit,
%point meta=explicit symbolic,
]
table[
y=actualClass,x=estimatedClass,meta=siftMean
 ]{lowLevelFeaturesConfusion.dat};
 
\node[small dot,pin=120:{xlabel anchor}]     at (axis description cs:0.5,1.1) {};           
\node[small dot,pin=-30:{ylabel anchor}]     at (axis description cs:-0.1,0.5) {};         

\end{axis}

\begin{axis}[gtLabelsAxis,
			 at=(siftConfMatrix.north),anchor=south,
			 yshift=\gtLabelDistance,
			 name=xGTLabels,
			 ]
			\addplot [gtLabelsPlotStyle,point meta=x] {0};
\end{axis}

 
\begin{axis}[gtLabelsAxis,%rotate=-90,yshift=-\gtLabelDistance,
			height=0cm,
			 name=yGTLabels,clip=false,
			 ]
			\addplot [gtLabelsPlotStyle,point meta=x] {0};   
%			\node[small dot,pin=125:{putaaaaaa}] at (axis description cs:0.5,0.5) {}; 
%			\node[small dot,pin=120:{$(0,0)$}]     at (axis description cs:0,0) {};           
%\node[small dot,pin=-30:{$(1,1)$}]     at (axis description cs:1,1) {};         
%\node[small dot,pin=-90:{$(1.03,0.5)$}]at (axis description cs:1.03,0.5) {}; 

			
\end{axis}


\end{tikzpicture}

\begin{tikzpicture}

\end{tikzpicture}



\end{document}
  \label{fig:llConfMatrix}
\end{tikzfigure}
} 

\end{columns}

\block{\ac{sift} as a high-level descriptor using \acf{bof}, tested using \ac{rbf}-\ac{svm} classifier}{
  \begin{multicols*}{3}
  %Texture is an area property related to spatial repetition of structures, statistical similarities, or both.

  \begin{itemize}

    \item Goal: to analyze tissue classes separability based on
      \begin{itemize}
        \item Classification framework for tissue discrimination (\ac{rbf}-\ac{svm}).
        \item High-level descriptors based on \ac{bof}-\ac{sift}.
      \end{itemize}

    \item Feature Description (see fig.\,\ref{fig:bof}).
      \begin{itemize}
        \item Generate a codebook of the features (k-means, $k=36$).
        \item Extract superpixels (\ac{qs}).
        \item Occurrence study to describe each superpixel.
      \end{itemize}
      \vspace{1pt}
    \item Quantitative analysis can be found in fig.\,\ref{fig:hlConfMatrix},
          The sensitivity achieved is $29\pm3.6\%$ for the intensity and $33.5\pm2.3\%$ for \ac{sift}.
  \end{itemize}

%  In order to produce quantitative results, dataset of superpixels is generated based on the images and the multi-label \ac{gt} available.
%  For a superpixel to be eligible, an area larger than $75\%$ need to belong to a single \ac{gt} label. The resulting dataset contains 20 folds of 8 superpixels (one per class). 

%  The experiment has been repeated for different coodebooks to take into account the variability of the coodebook within the results. 
%  Figure~\ref{fig:hlConfMatrix} compares the results of using \ac{bof} based on \ac{sift} or intensity for encoding \ac{us} texture. 
%  The sensitivity achieved is $29\pm3.6\%$ for the intenisty and $33.5\pm2.3\%$ for \ac{sift}.
  \end{multicols*}
}

\begin{columns}
\column{0.60} \block{}{
\begin{tikzfigure}[ \acs{sift}-\acs{bof} descriptors qualitative analysis.
 Image example; dictionary example; dictionary occurrence associated with the highlighted superpixels.]
\documentclass[border=2pt]{standalone}
\usepackage{tikz}


\input{../../defineBreastGTlabelColors.tex}

\begin{document}
\begin{tikzpicture}
\node[anchor=south west,inner sep=0] (imgNode) at (0,0) {\includegraphics[trim = 91 300 100 242, clip,width=.2\textwidth]{appearance1.pdf}};

\tikzset{nameSt/.style=
{anchor=north west,rectangle,
node distance=10pt,
minimum width=8pt,minimum height=8pt,
inner sep=0,
}}
\begin{tiny}

\draw[] (imgNode.north east) +(5pt,0) node[nameSt, label=below:Codebook] (dictionary) {\includegraphics[width=.10\textwidth,]{dictionary/01.png}} ;
\draw[] (dictionary.north east) +(5pt,0) node[nameSt, label=below:(1)] (BoFi) {\includegraphics[width=.10\textwidth,]{dictionary/01Signatures/127.png}} ;
\draw[] (BoFi.north east) +(5pt,0) node[nameSt, label=below:(2)] (BoFii) {\includegraphics[width=.02\textwidth,]{dictionary/01Signatures/037.png}} ;

\node[nameSt, below=of dictionary, label=below:(3)] (BoFiii) {\includegraphics[width=.02\textwidth,]{dictionary/01Signatures/054.png}} ;
\node[nameSt, below=of BoFi, label=below:(4)] (BoFiv) {\includegraphics[width=.02\textwidth,]{dictionary/01Signatures/118.png}} ;
\node[nameSt, below=of BoFii, label=below:(5)] (BoFv) {\includegraphics[width=.02\textwidth,]{dictionary/01Signatures/160.png}} ;

\node[nameSt, below=of BoFiii, label=below:(6)] (BoFvi) {\includegraphics[width=.02\textwidth,]{dictionary/01Signatures/061.png}} ;
\node[nameSt, below=of BoFiv, label=below:(7)] (BoFvii) {\includegraphics[width=.02\textwidth,]{dictionary/01Signatures/077.png}} ;
\node[nameSt, below=of BoFv, label=below:(8)] (BoFviii) {\includegraphics[width=.02\textwidth,]{dictionary/01Signatures/043.png}} ;

\end{tiny}

\end{tikzpicture}
\end{document}
 
\label{fig:bof}
\end{tikzfigure}
}

\column{0.40} \block{}{
\begin{tikzfigure}[Confusion matrix showing quantitative results obtained from ($1\text{ samples} \times 8 \text{ classes} \times 3 \text{ codebooks}) \times 20\text{-fold}$ cross-validation.
  ]
\documentclass[border=2pt]{standalone}
\usepackage{tikz}
\usepackage{pgfplots}
\usepackage{scalefnt}

\input{../../defineBreastGTlabelColors.tex}
\begin{document}
\makeatletter
  \newcommand\tinyv{\@setfontsize\tinyv{7pt}{6}}
\makeatother

\newcommand\confMatrixSize{9cm}
\newcommand\gtLabelDistance{0.3cm}

\pgfplotsset{
	myStyleCommon/.append style={	width=\confMatrixSize,
                                height=\confMatrixSize,
                                scale only axis,
                                xlabel={\small Estimated Class},
                                ylabel={\small Actual Class},
                                yticklabels={},
                                xticklabels={},
                                xticklabel pos=right,
                                xlabel near ticks,
                                xtick={1.5,2.5,...,7.5},
                                ytick={1.5,2.5,...,7.5},
                                grid,
                              } }

\pgfplotsset{
	confMatrix/.append	style={	myStyleCommon,
                              axis on top,
                              colormap/jet,
                              point meta min=0,
                              point meta max=1,
                              xmin=0.5,
                              xmax=8.5,
                              y dir=reverse,
                              ymin=0.5,
                              ymax=8.5,
							              } }

\pgfplotsset{							  
	gtLabelsAxis/.append style= { myStyleCommon,
                                axis lines=none,
                                colormap={labelsGT}{color(0cm)=(bgColor);
                                          color(1cm)=(lungColor);	
                                          color(2cm)=(chestWallColor);
                                          color(3cm)=(ribColor);
                                          color(4cm)=(pectoralColor);
                                          color(5cm)=(fibroGlandColor); 
                                          color(6cm)=(fatColor); 
                                          color(7cm)=(skinColor);    
                                          color(8cm)=(lesionColor);
                                          color(9cm)=(boundaryColor); 
                                          color(10cm)=(unkTissueColor); 
                                          },
                                  point meta min=1,
                                  point meta max=8,
                                  xmin=-0.5,xmax=7.5,
                                  height=0cm,
							                } }

\tikzstyle{gtLabelsPlotStyle}= [scatter,
                                only marks,
                                mark=square*,
                                mark size=5pt,
                                domain=1:8,
                                samples at={1,...,8},
                               ]

\tikzstyle{stdDrawingStyle} = [scatter, only marks, no marks,]
\tikzstyle{meanValueStyle}  = [ font=\tinyv,
                                fill=white,
                                inner sep = 2pt,
                                fill opacity=0.5,
                                text opacity=1,
                                %anchor=\coordindex,
                              ]
\tikzstyle{stdValuesStyle}  = [ meanValueStyle,
                                anchor=north,
                              ]
 

\begin{tikzpicture}
\begin{axis}
  [ confMatrix,
    %colorbar,colorbar to name={storedcolorbar},
    name=siftConfMatrix,
    %title=SIFT,
  ]
  \addplot graphics [xmin=0.5,xmax=8.5,ymin=0.5,ymax=8.5] 
  {/home/sik/Work/escola/recerca/iwdm2014/poster/figures/highLevelConfMatrix/highLevelSIFTConfusionMatrixMean.png};

  \addplot
    [ stdDrawingStyle,
      % we use ’point meta’ as color data...
      point meta=\thisrow{siftStd},
      % ... therefore, we can’t use it as argument for nodes near coords ..
      % ... look how to print the numbers at pgfplotstable.pdf manual
      nodes near coords*=
        {$
        \pgfmathprintnumber[fixed,precision=0]{\meanVal}\%
        $},
      % ... which requires to define a visualization dependency:
      visualization depends on={100*\thisrow{siftMean} \as \meanVal},
      % Add the node visual cues
      every node near coord/.append style= {meanValueStyle },
    ]
  table[ y=actualClass, x=estimatedClass ]
  {/home/sik/Work/escola/recerca/iwdm2014/poster/figures/highLevelConfMatrix/highLevelFeaturesConfusion.dat};

  \addplot % Like the previous plot, but this time with less unncecessary stuff.
    [ stdDrawingStyle,
      nodes near coords*=
        {$
        \pm\hspace{-0.1em}
        \pgfmathprintnumber[fixed,precision=1]{\stdVal}
        $},
      visualization depends on={100*\thisrow{siftStd} \as \stdVal},
      every node near coord/.append style= {stdValuesStyle},
    ]
  table[ y=actualClass, x=estimatedClass ]
  {/home/sik/Work/escola/recerca/iwdm2014/poster/figures/highLevelConfMatrix/highLevelFeaturesConfusion.dat};

\end{axis}

\begin{axis}[gtLabelsAxis,
			 at=(siftConfMatrix.north),anchor=south,
			 yshift=\gtLabelDistance,
			 name=xGTLabelsSIFT,
			 ]
			\addplot [gtLabelsPlotStyle,point meta=x] {0};
\end{axis}
 
\begin{axis}[	gtLabelsAxis,
				rotate=-90,yshift=-\gtLabelDistance,
			 	name=yGTLabelsSIFT,
			 ]
			\addplot [gtLabelsPlotStyle,point meta=x] {0};   	
\end{axis}
\node at (siftConfMatrix.south) 
      [ anchor=north,
        yshift=-15pt,
      ]{SIFT};
\end{tikzpicture}
%intensityMean	intensityStd
%
\begin{tikzpicture}
\begin{axis}
  [ confMatrix,
    %colorbar,colorbar to name={storedcolorbar},
    name=intensityConfMatrix,
    %title=SIFT,
  ]
  \addplot graphics [xmin=0.5,xmax=8.5,ymin=0.5,ymax=8.5] 
  {/home/sik/Work/escola/recerca/iwdm2014/poster/figures/highLevelConfMatrix/highLevelIntensityConfusionMatrixMean.png};

  \addplot
    [ stdDrawingStyle,
      % we use ’point meta’ as color data...
      point meta=\thisrow{siftStd},
      % ... therefore, we can’t use it as argument for nodes near coords ..
      % ... look how to print the numbers at pgfplotstable.pdf manual
      nodes near coords*=
        {$
        \pgfmathprintnumber[fixed,precision=0]{\meanVal}\%
        $},
      % ... which requires to define a visualization dependency:
      visualization depends on={100*\thisrow{intensityMean} \as \meanVal},
      % Add the node visual cues
      every node near coord/.append style= {meanValueStyle },
    ]
  table[ y=actualClass, x=estimatedClass ]
  {/home/sik/Work/escola/recerca/iwdm2014/poster/figures/highLevelConfMatrix/highLevelFeaturesConfusion.dat};

  \addplot % Like the previous plot, but this time with less unncecessary stuff.
    [ stdDrawingStyle,
      nodes near coords*=
        {$
        \pm\hspace{-0.1em}
        \pgfmathprintnumber[fixed,precision=1]{\stdVal}
        $},
      visualization depends on={100*\thisrow{intensityStd} \as \stdVal},
      every node near coord/.append style= {stdValuesStyle},
    ]
  table[ y=actualClass, x=estimatedClass ]
  {/home/sik/Work/escola/recerca/iwdm2014/poster/figures/highLevelConfMatrix/highLevelFeaturesConfusion.dat};

\end{axis}

\begin{axis}[gtLabelsAxis,
			 at=(siftConfMatrix.north),anchor=south,
			 yshift=\gtLabelDistance,
			 name=xGTLabelsIntensity,
			 ]
			\addplot [gtLabelsPlotStyle,point meta=x] {0};
\end{axis}
 
\begin{axis}[	gtLabelsAxis,
				rotate=-90,yshift=-\gtLabelDistance,
			 	name=yGTLabelsIntensity,
			 ]
			\addplot [gtLabelsPlotStyle,point meta=x] {0};   	
\end{axis}
\node at (intensityConfMatrix.south) 
      [ anchor=north,
        yshift=-10pt,
      ]{Intensity};
\end{tikzpicture}
\end{document}

\label{fig:hlConfMatrix}
\end{tikzfigure}
} 
\end{columns}

%\block[titleoffsety=-1cm,bodyoffsety=-1cm]{Sample document}{This poster...}
\block{Conclusion}{
The present study was designed to explore the usage of \ac{sift} feature space as a texture for characterizing the different tissues present in a breast \ac{us} image. 
The usage of \ac{sift} either as a low-level or high-level texture descriptor has been evaluated in comparison to intensity features, which are the features most commonly used.
The fact that \ac{sift} and intensity descriptors produce similar results, encourages further studies on using \ac{sift} texture descriptors characterizing breast tissues in \ac{us} images.

}
%\note[targetoffsetx=24cm, targetoffsety=-9cm,radius=8cm,width=.75\textwidth,innersep=.4cm]{You can...}
\note[targetoffsetx=-8cm,targetoffsety=-5cm,radius=8cm,width=.8\textwidth,innersep=.4cm]{\centering This work was partially supported by the Spanish Science and Innovation grants nb.\,TIN2012-37171-C02-01 and TIN2012-37171-C02-02, and the Regional Council of Burgundy.}
\end{document}
